% Festlegung des Allgemeinen Dokumentenformats
\documentclass[a4paper,12pt,headsepline]{scrartcl}

% Umlaute unter UTF8 nutzen
\usepackage[utf8]{inputenc}

% Variablen
%Variablen welche innerhalb der gesamten Arbeit zur Verfügung stehen sollen
\newcommand{\titleDocument}{Bachelorarbeit}
\newcommand{\subjectDocument}{
    Analyse der Echtzeitfähigkeit von Micro-ROS und FreeRTOS am Beispiel einer
    Robotersteuerungssoftware
}


% weitere Pakete
% Grafiken aus PNG Dateien einbinden
\usepackage{graphicx}
\usepackage{tikz}

% Deutsche Sonderzeichen und Silbentrennung nutzen
\usepackage[ngerman]{babel}

% Eurozeichen einbinden
\usepackage[right]{eurosym}

% Zeichenencoding
\usepackage[T1]{fontenc}

\usepackage{lmodern}

% floatende Bilder ermöglichen
%\usepackage{floatflt}

% mehrseitige Tabellen ermöglichen
\usepackage{longtable}

% Unterstützung für Schriftarten
%\newcommand{\changefont}[3]{ 
%\fontfamily{#1} \fontseries{#2} \fontshape{#3} \selectfont}

% Packet für Seitenrandabständex und Einstellung für Seitenränder
\usepackage{geometry}
\geometry{left=3.5cm, right=2cm, top=2.5cm, bottom=2cm}

% Paket für Boxen im Text
\usepackage{fancybox}

% bricht lange URLs "schön" um
\usepackage[hyphens,obeyspaces,spaces]{url}

% Paket für Textfarben
\usepackage{color}

% Mathematische Symbole importieren
\usepackage{amssymb}

\usepackage{amsmath}

% auf jeder Seite eine Überschrift (alt, zentriert)
%\pagestyle{headings}

\usepackage{pdfpages}

% erzeugt Inhaltsverzeichnis mit Querverweisen zu den Abschnitten (PDF Version)
\usepackage[bookmarksnumbered,pdftitle={\titleDocument},hyperfootnotes=false]{hyperref}
%\hypersetup{colorlinks, citecolor=red, linkcolor=blue, urlcolor=black}
%\hypersetup{colorlinks, citecolor=black, linkcolor= black, urlcolor=black}

% neue Kopfzeilen mit fancypaket
\usepackage{fancyhdr} %Paket laden
\pagestyle{fancy} %eigener Seitenstil
\fancyhf{} %alle Kopf- und Fußzeilenfelder bereinigen
\fancyhead[L]{\nouppercase{\leftmark}} %Kopfzeile links
\fancyhead[C]{} %zentrierte Kopfzeile
\fancyhead[R]{\thepage} %Kopfzeile rechts
\renewcommand{\headrulewidth}{0.4pt} %obere Trennlinie
%\fancyfoot[C]{\thepage} %Seitennummer
%\renewcommand{\footrulewidth}{0.4pt} %untere Trennlinie

% change font to serif for headings - buggy
% \usepackage{sectsty}
% \allsectionsfont{\normalfont\bfseries}

% für Tabellen
\usepackage{array}

% Runde Klammern für Zitate
%\usepackage[numbers,round]{natbib}

% Festlegung Art der Zitierung - Havardmethode: Abkuerzung Autor + Jahr
\bibliographystyle{alphadin}

% Schaltet den zusätzlichen Zwischenraum ab, den LaTeX normalerweise nach einem Satzzeichen einfügt.
%\frenchspacing

% Paket für Zeilenabstand
\usepackage{setspace}

% für Bildbezeichner
\usepackage{capt-of}

% für Stichwortverzeichnis
\usepackage{makeidx}

%Konfiguriere das Inhaltsverzeichnis
\usepackage{tocbasic}
\DeclareTOCStyleEntries[
  raggedentrytext,
  numwidth=0pt,
  numsep=1ex,
  dynnumwidth,
]{tocline}{chapter,section,subsection,subsubsection,paragraph,subparagraph}
\DeclareTOCStyleEntries[
  indent=0pt,
  linefill=\TOCLineLeaderFill,
]{tocline}{section,subsection,subsubsection,paragraph,subparagraph}


% für Listings
% \usepackage{listings}
% \lstset{numbers=left, numberstyle=\tiny, numbersep=5pt, keywordstyle=\color{black}\bfseries, stringstyle=\ttfamily,showstringspaces=false,basicstyle=\footnotesize,captionpos=b}
\usepackage{xcolor}
\usepackage[newfloat]{minted}
\usepackage{caption}
\usepackage{fancyhdr}
\newenvironment{code}{\captionsetup{type=listing}}{}
\SetupFloatingEnvironment{listing}{name=Quellcode}
\setminted{
    linenos,
	frame=single,
    bgcolor=black!5,
    fontsize=\footnotesize
}
\usepackage{graphicx}  % Add this to your preamble

\usepackage[withpage]{acronym}

% Absatz ohne Einrückung
\setlength{\parskip}{1em} % 1em entspricht der Breite eines 'M'-Zeichens
\setlength{\parindent}{0pt}

% Indexerstellung
\makeindex

% Abkürzungsverzeichnis
\usepackage[german]{nomencl}
\let\abbrev\nomenclature

% Abkürzungsverzeichnis LiveTex Version
% Titel des Abkürzungsverzeichnisses
\renewcommand{\nomname}{Abkürzungsverzeichnis}
% Abstand zwischen Abkürzung und Erläuterung
\setlength{\nomlabelwidth}{.25\textwidth}
% Zwischenraum zwischen Abkürzung und Erläuterung mit Punkten
\renewcommand{\nomlabel}[1]{#1 \dotfill}
% Variation des Abstandes der einzelnen Abkürzungen zu einander
\setlength{\nomitemsep}{-\parsep}
% Index mit Abkürzungen erzeugen
\makenomenclature
%\makeglossary

% Abkürzungsverzeichnis TeTEX Version
% \usepackage[german]{nomencl}
% \makenomenclature
% %\makeglossary
% \renewcommand{\nomname}{Abkürzungsverzeichnis}
% \AtBeginDocument{\setlength{\nomlabelwidth}{.25\columnwidth}}
% \renewcommand{\nomlabel}[1]{#1 \dotfill}
% \setlength{\nomitemsep}{-\parsep}

% Optional: Einzelne Zeilen am Anfang einer Seite unterdrücken (Schusterjungen)
% \clubpenalty = 10000
% Optional: Einzelne Zeilen am Ende einer Seite unterdrücken (Hurenkinder)
% \widowpenalty = 10000
% \displaywidowpenalty = 10000

\begin{document}
% hier werden die Trennvorschläge inkludiert
%hier müssen alle Wörter rein, welche Latex von sich auch nicht korrekt trennt bzw. bei denen man die genaue Trennung vorgeben möchte
\hyphenation{
Film-pro-du-zen-ten
Lux-em-burg
Soft-ware-bau-steins
zeit-in-ten-siv
}


% Schriftart Helvetica verwenden
%\usepackage{helvet}
%\renewcommand\familydefault{\sfdefault}

% Leere Seite am Anfang
%\thispagestyle{empty} % erzeugt Seite ohne Kopf- / Fusszeile
%\mbox{}
%\newpage

% Titelseite %
\thispagestyle{empty}

% % Logo
% \begin{figure}[t]
%  \centering
%  \includegraphics[width=0.3\textwidth]{assets/fhlogo.pdf}
% \end{figure}

\begin{center}
\textbf{\Large{~\titleDocument}}
\end{center}
\begin{verbatim}





\end{verbatim}
\begin{center}
\doublespacing
\textbf{\LARGE{\subjectDocument}}\\
\singlespacing
\end{center}
\begin{verbatim}
















\end{verbatim}
\begin{center}
\textbf{
    An der Fachhochschule Dortmund \\
    im Fachbereich Informatik \\
    Studiengang Technische Informatik \\
    erstellte Thesis \\
    zur Erlangung des akademischen Grades \\
    Bachelor of Science \\
    B. Sc. \\
}
\end{center}
\begin{center}
    \textbf{Xu, Zijian \linebreak
    geboren am 25.09.1998 \linebreak
    7204211
    }
\end{center}
\begin{flushleft}
\begin{tabular}{llll}
\textbf{Betreuung durch:} & & Prof. Dr. Christof Röhrig &\\
    & & M. Sc. Alexander Miller &\\
\textbf{Version vom:} & & Dortmund, \today &\\
\end{tabular}
\end{flushleft}


% römische Numerierung
\pagenumbering{roman}

% 1.5 facher Zeilenabstand
\onehalfspacing

\newpage

% Einleitung / Abstract
\thispagestyle{empty}
\section*{Kurzfassung}

Diese Arbeit analysiert die Echtzeitfähigkeit von Micro-ROS und FreeRTOS anhand
einer Robotersteuerungssoftware. Ziel ist der Vergleich beider Systeme
hinsichtlich ihres Echtzeitverhaltens mit Schwerpunkt auf den Ausführungszeiten.
Dazu wird zunächst die bestehende Micro-ROS-Steuerungssoftware auf FreeRTOS
portiert. Anschließend wird eine zyklengenaue Messmethode für den Programmlauf
basierend auf der Data Watchpoint and Trace Unit (DWT) der ARM-Architektur
entwickelt. Zum Abschluss werden die Laufzeitdaten visualisiert und ausgewertet.
Die Ergebnisse verdeutlichen die unterschiedlichen Schwerpunkte beider
Plattformen: Micro-ROS punktet durch die Kompatibilität mit dem ROS-Ökosystem,
während FreeRTOS mit minimalen Latenzen und deterministischem Scheduling
deutlich besseres Echtzeitverhalten bietet. Die Analyse bestätigt zudem den
maßgeblichen Einfluss der Cache-Nutzung auf die Rechenleistung.

\section*{Abstract}

This thesis analyzes the real-time capabilities of micro-ROS and FreeRTOS using
a robotic control software. The goal is to compare both systems regarding their
real-time behavior emphasized on the execution times. First, the existing
micro-ROS control software is ported to FreeRTOS. Next, a cycle-accurate
measurement method for program execution based on ARM's Data Watchpoint and
Trace Unit (DWT)is developed. Finally, the runtime data is evaluated. The
evaluation of the results highlights the different strengths of both platforms:
micro-ROS stands out due to its compatibility with the ROS ecosystem, while
FreeRTOS achieves superior real-time behavior through minimal latencies and
deterministic scheduling. Additionally, the analysis confirms the significant
impact of the cache utilization on system performance.


% einfacher Zeilenabstand
\singlespacing

\newpage
% Seitenzählung bei Inhaltsverzeichnis beginnen
\setcounter{page}{1}

% Inhaltsverzeichnis anzeigen
\thispagestyle{empty}
\tableofcontents

\newpage
% das Abbildungsverzeichnis
% Verion 1: Abbildungsverzeichnis MIT führender Nummberierung endgueltig anzeigen
\listoffigures
% Abbildungsverzeichnis soll im Inhaltsverzeichnis auftauchen
\addcontentsline{toc}{section}{Abbildungsverzeichnis}

% Verion 2: Abbildungsverzeichnis OHNE führende Nummberierung endgueltig anzeigen
%\begingroup
%\renewcommand\numberline[1]{}
%\listoffigures
%\endgroup


% % das Tabellenverzeichnis
% \newpage
% % \fancyhead[L]{Abbildungsverzeichnis / Abkürzungsverzeichnis} %Kopfzeile links
% % Tabellenverzeichnis endgültig anzeigen
% \listoftables
% % Tabellenverzeichnis soll im Inhaltsverzeichnis auftauchen
% \addcontentsline{toc}{section}{Tabellenverzeichnis}

% das Quellcodeverzeichnis
\newpage
\renewcommand*{\listlistingname}{Quellcodeverzeichnis}
\listoflistings % Add Quellcodeverzeichnis
\addcontentsline{toc}{section}{Quellcodeverzeichnis}

% das Abkürzungsverzeichnis
\newpage
% das Abkürzungsverzeichnis ausgeben
\fancyhead[L]{Abkürzungsverzeichnis} %Kopfzeile links
\section*{Abkürzungsverzeichnis}

\begin{acronym}
\end{acronym}

% \printnomenclature[3cm]
% Abkürzungsverzeichnis soll im Inhaltsverzeichnis auftauchen
\addcontentsline{toc}{section}{Abkürzungsverzeichnis}


%%%%%%% EINLEITUNG %%%%%%%%%%%%
\newpage
\fancyhead[L]{\nouppercase{\leftmark}} %Kopfzeile links

% 1,5 facher Zeilenabstand
\onehalfspacing

% arabische Seitennummerierung ab hier
\pagenumbering{arabic}

% Alternative Einbindung des Abstract in Kapitel "0" falls gewünscht
%\setcounter{section}{-1}
%\setcounter{page}{0}

% Option: Einbindung abstract
%\section*{Kurzfassung}

Diese Arbeit analysiert die Echtzeitfähigkeit von Micro-ROS und FreeRTOS anhand
einer Robotersteuerungssoftware. Ziel ist der Vergleich beider Systeme
hinsichtlich ihres Echtzeitverhaltens mit Schwerpunkt auf den Ausführungszeiten.
Dazu wird zunächst die bestehende Micro-ROS-Steuerungssoftware auf FreeRTOS
portiert. Anschließend wird eine zyklengenaue Messmethode für den Programmlauf
basierend auf der Data Watchpoint and Trace Unit (DWT) der ARM-Architektur
entwickelt. Zum Abschluss werden die Laufzeitdaten visualisiert und ausgewertet.
Die Ergebnisse verdeutlichen die unterschiedlichen Schwerpunkte beider
Plattformen: Micro-ROS punktet durch die Kompatibilität mit dem ROS-Ökosystem,
während FreeRTOS mit minimalen Latenzen und deterministischem Scheduling
deutlich besseres Echtzeitverhalten bietet. Die Analyse bestätigt zudem den
maßgeblichen Einfluss der Cache-Nutzung auf die Rechenleistung.

\section*{Abstract}

This thesis analyzes the real-time capabilities of micro-ROS and FreeRTOS using
a robotic control software. The goal is to compare both systems regarding their
real-time behavior emphasized on the execution times. First, the existing
micro-ROS control software is ported to FreeRTOS. Next, a cycle-accurate
measurement method for program execution based on ARM's Data Watchpoint and
Trace Unit (DWT)is developed. Finally, the runtime data is evaluated. The
evaluation of the results highlights the different strengths of both platforms:
micro-ROS stands out due to its compatibility with the ROS ecosystem, while
FreeRTOS achieves superior real-time behavior through minimal latencies and
deterministic scheduling. Additionally, the analysis confirms the significant
impact of the cache utilization on system performance.

%\newpage

% einzelne Kapitel werden hier eingebunden
\section*{Einleitung}

Die vorliegende Arbeit beschäftigt sich mit der Analyse der Echtzeitfähigkeit
von Micro-ROS und FreeRTOS am Beispiel einer Robotersteuerungssoftware. Ziel ist
es, die Performance beider Systeme hinsichtlich der Ausführungszeiten,
Ressourcenverbrauch sowie Echtzeitverhalten zu untersuchen, um ihre Eignung für
Roboteranwendungen zu bewerten.

Die Arbeit beinhaltet schwerpunktmäßig die Entwicklung einer Methode zum
Profiling der Steuerungssoftware eines mobilen Roboters. Dabei wird zunächst die
bestehende Firmware, die auf Micro-ROS basiert, im Rahmen dieser Arbeit auf
FreeRTOS portiert. Anschließend werden die Methodik zur Generierung von
Profiling-Daten für die Analyse detailliert beschrieben und die resultierenden
Ergebnisse evaluiert.

Zu Beginn wird ein Überblick über die grundlegenden Konzepte gegeben.
Darauffolgend werden die Implementierungen detailliert vorgestellt. Abschließend
werden die erzielten Ergebnisse evaluiert, und es wird ein Ausblick auf weitere
Anwendungsmöglichkeiten und Optimierungspotenziale gegeben.

\newpage

% TODO

\section{Abschluss}

Am Anfang wurde die Robotersteuerungssoftware von Micro-ROS auf FreeRTOS
umgestellt. Anschließend wurden eine Multi-Producer-Senke sowie ein Verfahren
entwickelt, das Echtzeitinformationen über die Steuerungssoftware ausgeben kann.
Abschließend wurde die Echtzeitfähigkeit basierend auf die Echtzeitinformationen
analysiert.

\subsection{Fazit}

Es lässt sich schlussfolgern, dass die Steuerungssoftware zwar durch Integration
von Micro-ROS funktionsreicher und folglich mit einer Vielzahl von
ROS-Komponenten kompatibel wird, dies allerdings mit erheblichem Overhead
erkauft wird. Bei begrenztem Speicher oder Rechenleistung bleibt FreeRTOS mit
seinem schlanken Kernel und den threadsicheren Queue-Implementierungen weiterhin
eine geeignete Wahl gegenüber komplexeren RTOS-Lösungen -- besonders wenn harte
Echtzeitfähigkeit im Vordergrund steht.

Darüber hinaus wurde gezeigt, dass sich die Softwareleistung durch L1-Caches
deutlich verbessern lässt, was für leistungskritische Software entscheidend ist.

\subsection{Ausblick}

Für zukünftige Arbeiten könnte die Multi-Producer-Senke so weiterentwickelt
werden, dass sie atomare Schreiboperationen auf 4-Byte-/32-Bit-Ebene
unterstützt. Dadurch könnten die Echtzeit-Daten nicht mehr im menschenlesbaren
Format, sondern maschinenlesbar und komprimiert als 4-Byte-Dateneinheit
ausgegeben werden.

Dies würde den Zwischenpuffer für die Echtzeitanalyse überflüssig machen, da die
durch Kontextwelchsel/Interrupt generierten Daten direkt atomar in die Senke
geschrieben werden könnten. Dazu müsste ein Parser auf dem Host-System
entwickelt werden, der idealerweise auch die Visualisierung sowie Analyse der
Daten parallel in Echtzeit ermöglicht.

\newpage

% einfacher Zeilenabstand
\singlespacing
% Literaturliste soll im Inhaltsverzeichnis auftauchen
\newpage
\phantomsection
\addcontentsline{toc}{section}{Literaturverzeichnis}
% Literaturverzeichnis anzeigen
\renewcommand\refname{Literaturverzeichnis}
\bibliography{Hauptdatei}

%% Index soll Stichwortverzeichnis heissen
% \newpage
% % Stichwortverzeichnis soll im Inhaltsverzeichnis auftauchen
% \addcontentsline{toc}{section}{Stichwortverzeichnis}
% \renewcommand{\indexname}{Stichwortverzeichnis}
% % Stichwortverzeichnis endgültig anzeigen
% \printindex

% \onehalfspacing
% % evtl. Anhang
% \newpage
% \phantomsection
% \addcontentsline{toc}{section}{Anhang}
% \fancyhead[L]{Anhang} %Kopfzeile links
% \input{anhang}

\newpage
% \includepdf[pages=1, fitpaper=true]{eigenstaendigkeitserklaehrung.pdf}

% \fancyhead[L]{Eigenständigkeitserklährung} %Kopfzeile links
% \input{eigenstaendigkeitserklaehrung} % deleted 

\end{document}
